% Part: turing-machines
% Chapter: undecidability
% Section: verification

\documentclass[../../../include/open-logic-section]{subfiles}

\begin{document}

\olfileid{tur}{und}{ver}
\olsection{التحقق من صحة التمثيل}

\begin{explain}
صحة هذا التمثيل تتم بإثبات جهتين. الأولى: إذا توقفت $\OLId{latin-looped}{M}$ عند
الدخل~$\OLId{latin-ordinary}{w}$، صدقت $!T(\OLId{latin-looped}{M},\OLId{latin-ordinary}{w})\lif !E(\OLId{latin-looped}{M},\OLId{latin-ordinary}{w})$ منطقيًا. والثانية عكسها:
إذا كانت $!T(\OLId{latin-looped}{M},\OLId{latin-ordinary}{w})\lif !E(\OLId{latin-looped}{M},\OLId{latin-ordinary}{w})$ صادقة منطقيًا، لزم أن تتوقف
$\OLId{latin-looped}{M}$ فعلًا في نهاية المطاف عند الدخل~$\OLId{latin-ordinary}{w}$.

ولكل جهة طريق غير طريق الأخرى. ففي الأولى نريد إثبات الصدق
المنطقي ل!!a{sentence} من الرتبة الأولى، هي
$!T(\OLId{latin-looped}{M},\OLId{latin-ordinary}{w})\lif !E(\OLId{latin-looped}{M},\OLId{latin-ordinary}{w})$. وأيسر ذلك أن نقيم لها !!a{derivation}.
غير أن الحكم عام في $\OLId{latin-looped}{M}$ و$\OLId{latin-ordinary}{w}$؛ فليس المطلوب !!{sentence}
واحدة نقيم لها !!a{derivation}، بل جمل لا تتناهى. فلا بد من
حجة استقرائية تبيّن أنه، أيًا كانت $\OLId{latin-looped}{M}$ و~$\OLId{latin-ordinary}{w}$، وأيًا كان عدد
خطوات $\OLId{latin-looped}{M}$ قبل توقفها عند الدخل~$\OLId{latin-ordinary}{w}$، يوجد !!a{derivation}
لـ$!T(\OLId{latin-looped}{M},\OLId{latin-ordinary}{w})\lif !E(\OLId{latin-looped}{M},\OLId{latin-ordinary}{w})$.

ومحل الاستقراء عدد الخطوات التي تقطعها $\OLId{latin-looped}{M}$ إلى تهيئة التوقف.
فنثبت لكل خطوة~$\OLId{latin-ordinary}{n}$ من تشغيل $\OLId{latin-looped}{M}$ عند الدخل~$\OLId{latin-ordinary}{w}$، إلى أول توقف،
أن $!T(\OLId{latin-looped}{M},\OLId{latin-ordinary}{w})\Entails !C(\OLId{latin-looped}{M},\OLId{latin-ordinary}{w},\OLId{latin-ordinary}{n})$؛ حيث $!C(\OLId{latin-looped}{M},\OLId{latin-ordinary}{w},\OLId{latin-ordinary}{n})$ وصف صحيح
لتهيئة $\OLId{latin-looped}{M}$ عند~$\OLId{latin-ordinary}{w}$ بعد~$\OLId{latin-ordinary}{n}$ خطوات. فإذا توقفت $\OLId{latin-looped}{M}$ عند
الدخل~$\OLId{latin-ordinary}{w}$ بعد $\OLId{latin-ordinary}{n}$ خطوة، كانت $!C(\OLId{latin-looped}{M},\OLId{latin-ordinary}{w},\OLId{latin-ordinary}{n})$ وصفًا لتهيئة توقف.
ونثبت أيضًا أن $!C(\OLId{latin-looped}{M},\OLId{latin-ordinary}{w},\OLId{latin-ordinary}{n})\Entails !E(\OLId{latin-looped}{M},\OLId{latin-ordinary}{w})$ متى كانت
$!C(\OLId{latin-looped}{M},\OLId{latin-ordinary}{w},\OLId{latin-ordinary}{n})$ كذلك. فإذا توقفت $\OLId{latin-looped}{M}$ عند الدخل~$\OLId{latin-ordinary}{w}$، وجد~$\OLId{latin-ordinary}{n}$
تكون عنده $\OLId{latin-looped}{M}$ في تهيئة توقف. وحينئذ
$!T(\OLId{latin-looped}{M},\OLId{latin-ordinary}{w})\Entails !C(\OLId{latin-looped}{M},\OLId{latin-ordinary}{w},\OLId{latin-ordinary}{n})$، وتصف $!C(\OLId{latin-looped}{M},\OLId{latin-ordinary}{w},\OLId{latin-ordinary}{n})$ ذلك التوقف؛
فينضم إلى هذا $!C(\OLId{latin-looped}{M},\OLId{latin-ordinary}{w},\OLId{latin-ordinary}{n})\Entails !E(\OLId{latin-looped}{M},\OLId{latin-ordinary}{w})$، وينتج
$!T(\OLId{latin-looped}{M},\OLId{latin-ordinary}{w})\Entails !E(\OLId{latin-looped}{M},\OLId{latin-ordinary}{w})$، أي
$\Entails !T(\OLId{latin-looped}{M},\OLId{latin-ordinary}{w})\lif !E(\OLId{latin-looped}{M},\OLId{latin-ordinary}{w})$.

وأما العكس فنبدأ فيه بفرض $\Entails !T(\OLId{latin-looped}{M},\OLId{latin-ordinary}{w})\lif !E(\OLId{latin-looped}{M},\OLId{latin-ordinary}{w})$،
ونطلب إثبات توقف $\OLId{latin-looped}{M}$ عند الدخل~$\OLId{latin-ordinary}{w}$. والفرض يفيد
$!T(\OLId{latin-looped}{M},\OLId{latin-ordinary}{w})\Entails !E(\OLId{latin-looped}{M},\OLId{latin-ordinary}{w})$؛ أي تصدق $!E(\OLId{latin-looped}{M},\OLId{latin-ordinary}{w})$ في كل
!!{structure} تصدق فيها $!T(\OLId{latin-looped}{M},\OLId{latin-ordinary}{w})$. فننشئ
!!a{structure}~$\Struct{M}$ تحقق $!T(\OLId{latin-looped}{M},\OLId{latin-ordinary}{w})$، مجالها $\Nat$؛ ومن تهيئات~$\OLId{latin-looped}{M}$ في تشغيلها عند الدخل~$\OLId{latin-ordinary}{w}$
نأخذ تفسيرات $\Obj Q_\OLId{latin-ordinary}{q}$ و$\Obj S_\sigma$. فيكون، مثلًا،
$\Sat{\OLId{latin-looped}{M}}{\Obj Q_\OLId{latin-ordinary}{q}(\num{\OLId{latin-ordinary}{m}},\num{\OLId{latin-ordinary}{n}})}$ إذا وفقط إذا كانت $\OLId{latin-looped}{M}$
عند الدخل~$\OLId{latin-ordinary}{w}$، بعد $\OLId{latin-ordinary}{n}$ خطوة، في الحالة~$\OLId{latin-ordinary}{q}$ ماسحةً الخانة~$\OLId{latin-ordinary}{m}$.
فلدينا $!T(\OLId{latin-looped}{M},\OLId{latin-ordinary}{w})\Entails !E(\OLId{latin-looped}{M},\OLId{latin-ordinary}{w})$ بالفرض،
و$\Sat{\OLId{latin-looped}{M}}{!T(\OLId{latin-looped}{M},\OLId{latin-ordinary}{w})}$ بالبناء؛ فيلزم $\Sat{\OLId{latin-looped}{M}}{!E(\OLId{latin-looped}{M},\OLId{latin-ordinary}{w})}$.
لكن $\Sat{\OLId{latin-looped}{M}}{!E(\OLId{latin-looped}{M},\OLId{latin-ordinary}{w})}$ تكافئ وجود $\OLId{latin-ordinary}{n}\in\Domain{\OLId{latin-looped}{M}}=\Nat$
تكون فيه $\OLId{latin-looped}{M}$ عند الدخل~$\OLId{latin-ordinary}{w}$ في تهيئة توقف بعد~$\OLId{latin-ordinary}{n}$ خطوات.
\end{explain}

\begin{defn}
نعرّف $!C(\OLId{latin-looped}{M},\OLId{latin-ordinary}{w},\OLId{latin-ordinary}{n})$ بأنها ال!!{sentence}
\[
\Obj Q_\OLId{latin-ordinary}{q}(\num{\OLId{latin-ordinary}{m}}, \num{\OLId{latin-ordinary}{n}}) \land \Obj S_{\sigma_\OLMathNumeral{0}}(\num{\OLMathNumeral{0}}, \num{\OLId{latin-ordinary}{n}})
\land \dots \land \Obj S_{\sigma_\OLId{latin-ordinary}{k}}(\num{\OLId{latin-ordinary}{k}}, \num{\OLId{latin-ordinary}{n}}) \land
\lforall[\OLId{latin-ordinary}{x}][(\num{\OLId{latin-ordinary}{k}} < \OLId{latin-ordinary}{x} \lif \Obj S_\TMblank(\OLId{latin-ordinary}{x}, \num{\OLId{latin-ordinary}{n}}))]
\]
وفيها $\OLId{latin-ordinary}{q}$ حالة $\OLId{latin-looped}{M}$ في الزمن~$\OLId{latin-ordinary}{n}$، وتكون $\OLId{latin-looped}{M}$ على الخانة~$\OLId{latin-ordinary}{m}$
في الزمن~$\OLId{latin-ordinary}{n}$؛ وتحمل الخانة~$\OLId{latin-ordinary}{i}$ الرمز~$\sigma_\OLId{latin-ordinary}{i}$ في الزمن~$\OLId{latin-ordinary}{n}$
لكل $\OLMathNumeral{0}\le \OLId{latin-ordinary}{i}\le \OLId{latin-ordinary}{k}$. وأما $\OLId{latin-ordinary}{k}$ فهو الأكبر من رقمين: رقم آخر خانة
غير فارغة إلى اليمين في الزمن~$\OLMathNumeral{0}$، ورقم أقصى خانة يمينًا زارها
الرأس خلال $\OLId{latin-ordinary}{n}$ خطوة.
\end{defn}

\begin{lem}\ollabel{lem:halt-config-implies-halt}
متى بلغت $\OLId{latin-looped}{M}$ عند الدخل~$\OLId{latin-ordinary}{w}$ تهيئة توقف بعد $\OLId{latin-ordinary}{n}$ خطوة، لزم
$!C(\OLId{latin-looped}{M},\OLId{latin-ordinary}{w},\OLId{latin-ordinary}{n})\Entails !E(\OLId{latin-looped}{M},\OLId{latin-ordinary}{w})$.
\end{lem}

\begin{proof}
لتتوقف $\OLId{latin-looped}{M}$ عند الدخل~$\OLId{latin-ordinary}{w}$ بعد $\OLId{latin-ordinary}{n}$ خطوة. فثمة حالة~$\OLId{latin-ordinary}{q}$ وخانة
$\OLId{latin-ordinary}{m}$ ورمز~$\sigma$ تحقق ما يأتي:
\begin{enumerate}
\item تكون $\OLId{latin-looped}{M}$ بعد $\OLId{latin-ordinary}{n}$ خطوة في الحالة~$\OLId{latin-ordinary}{q}$ على الخانة~$\OLId{latin-ordinary}{m}$
  التي تحمل~$\sigma$.
\item لا تكون دالة الانتقال $\delta(\OLId{latin-ordinary}{q},\sigma)$ معرّفة.
\end{enumerate}
ووصف التهيئة $!C(\OLId{latin-looped}{M},\OLId{latin-ordinary}{w},\OLId{latin-ordinary}{n})$ يتضمن
$\Obj Q_\OLId{latin-ordinary}{q}(\num{\OLId{latin-ordinary}{m}},\num{\OLId{latin-ordinary}{n}})$ و$\Obj S_\sigma(\num{\OLId{latin-ordinary}{m}},\num{\OLId{latin-ordinary}{n}})$.
ومن اجتماعهما تلزم $!E(\OLId{latin-looped}{M},\OLId{latin-ordinary}{w})$، أي:
\[
\lexists[\OLId{latin-ordinary}{x}][\lexists[\OLId{latin-ordinary}{y}][(\bigvee_{\tuple{\OLId{latin-ordinary}{q}, \sigma} \in
      \OLId{latin-looped}{X}}(\Obj Q_\OLId{latin-ordinary}{q}(\OLId{latin-ordinary}{x}, \OLId{latin-ordinary}{y}) \land \Obj S_\sigma(\OLId{latin-ordinary}{x}, \OLId{latin-ordinary}{y})))]]
\]
ووجه اللزوم أن $\Obj Q_\OLId{latin-ordinary}{q}(\num{\OLId{latin-ordinary}{m}},\num{\OLId{latin-ordinary}{n}})\land
\Obj S_\sigma(\num{\OLId{latin-ordinary}{m}},\num{\OLId{latin-ordinary}{n}})$ تقتضي
$\bigvee_{\tuple{\OLId{latin-ordinary}{q},\sigma}\in \OLId{latin-looped}{X}}(\Obj Q_\OLId{latin-ordinary}{q}(\num{\OLId{latin-ordinary}{m}},\num{\OLId{latin-ordinary}{n}})\land
\Obj S_\sigma(\num{\OLId{latin-ordinary}{m}},\num{\OLId{latin-ordinary}{n}}))$، وذلك لأن
$\tuple{\OLId{latin-ordinary}{q},\sigma}\in \OLId{latin-looped}{X}$.
\end{proof}

\begin{explain}
فمن توقف $\OLId{latin-looped}{M}$ عند الدخل $\OLId{latin-ordinary}{w}$ يلزم وجود~$\OLId{latin-ordinary}{n}$ بحيث
$!C(\OLId{latin-looped}{M},\OLId{latin-ordinary}{w},\OLId{latin-ordinary}{n})\Entails !E(\OLId{latin-looped}{M},\OLId{latin-ordinary}{w})$. وننتقل إلى إثبات أنه لكل زمن~$\OLId{latin-ordinary}{n}$
إلى أول تهيئة توقف يتحقق $!T(\OLId{latin-looped}{M},\OLId{latin-ordinary}{w})\Entails !C(\OLId{latin-looped}{M},\OLId{latin-ordinary}{w},\OLId{latin-ordinary}{n})$.
\end{explain}

\begin{lem}
\ollabel{lem:config}
لكل $\OLId{latin-ordinary}{n}$ لم تتوقف فيه $\OLId{latin-looped}{M}$ قبل تمام $\OLId{latin-ordinary}{n}$ خطوة، يكون
$!T(\OLId{latin-looped}{M},\OLId{latin-ordinary}{w})\Entails !C(\OLId{latin-looped}{M},\OLId{latin-ordinary}{w},\OLId{latin-ordinary}{n})$.
\end{lem}

\begin{proof}
أما أساس الاستقراء، ففي الحالة $\OLId{latin-ordinary}{n}=\OLMathNumeral{0}$ نجد كل مقترن من
$!C(\OLId{latin-looped}{M},\OLId{latin-ordinary}{w},\OLMathNumeral{0})$ مقترنًا في $!T(\OLId{latin-looped}{M},\OLId{latin-ordinary}{w})$ نفسها؛ فالاستلزام حاصل.

وأما فرض الاستقراء فهو: إن لم تتوقف $\OLId{latin-looped}{M}$ قبل الخطوة رقم $\OLId{latin-ordinary}{n}$،
كان $!T(\OLId{latin-looped}{M},\OLId{latin-ordinary}{w})\Entails !C(\OLId{latin-looped}{M},\OLId{latin-ordinary}{w},\OLId{latin-ordinary}{n})$. والمطلوب أن نبين، ما دامت
$!C(\OLId{latin-looped}{M},\OLId{latin-ordinary}{w},\OLId{latin-ordinary}{n})$ ليست وصف توقف، أن $!T(\OLId{latin-looped}{M},\OLId{latin-ordinary}{w})\Entails !C(\OLId{latin-looped}{M},\OLId{latin-ordinary}{w},\OLId{latin-ordinary}{n}+\OLMathNumeral{1})$.

لتكن $\OLId{latin-looped}{M}$ بعد $\OLId{latin-ordinary}{n}$ خطوة من ابتدائها عند $\OLId{latin-ordinary}{w}$ في الحالة~$\OLId{latin-ordinary}{q}$
على الخانة~$\OLId{latin-ordinary}{m}$. فإن $\OLId{latin-looped}{M}$، ما دامت لم تتوقف بعد~$\OLId{latin-ordinary}{n}$ خطوة،
لا بد أن يحتوي برنامج~$\OLId{latin-looped}{M}$ تعليمة من إحدى الصور الثلاث:

\begin{enumerate}
\item \ollabel{right} $\delta(\OLId{latin-ordinary}{q}, \sigma) = \tuple{\OLId{latin-ordinary}{q}', \sigma', \TMright}$

\item \ollabel{left} $\delta(\OLId{latin-ordinary}{q}, \sigma) = \tuple{\OLId{latin-ordinary}{q}', \sigma', \TMleft}$

\item \ollabel{stay} $\delta(\OLId{latin-ordinary}{q}, \sigma) = \tuple{\OLId{latin-ordinary}{q}', \sigma', \TMstay}$
\end{enumerate}

ولكل واحدة من هذه الصور حجة نوردها على الترتيب.

\begin{enumerate}
\item أما التعليمة من الصورة~\olref{right}، فبحسب
  \olref[rep]{defn:tm-descr}\olref[rep]{rep-right} تكون
\begin{align*}
& \lforall[\OLId{latin-ordinary}{x}][\lforall[\OLId{latin-ordinary}{y}][((\Obj Q_{\OLId{latin-ordinary}{q}}(\OLId{latin-ordinary}{x},
  \OLId{latin-ordinary}{y}) \land \Obj S_\sigma(\OLId{latin-ordinary}{x}, \OLId{latin-ordinary}{y})) \lif {}]]\\
& \qquad (\Obj Q_{\OLId{latin-ordinary}{q}'}(\OLId{latin-ordinary}{x}',
  \OLId{latin-ordinary}{y}') \land \Obj S_{\sigma'}(\OLId{latin-ordinary}{x}, \OLId{latin-ordinary}{y}') \land !A(\OLId{latin-ordinary}{x}, \OLId{latin-ordinary}{y})))
\intertext{مقترنة في $!T(M,w)$. وتستلزم ال!!{sentence} الآتية
  (بالتخصيص الكلي، مع $\num{m}$ بدل~$x$ و$\num{n}$ بدل~$y$):}
& (\Obj Q_{\OLId{latin-ordinary}{q}}(\num{\OLId{latin-ordinary}{m}}, \num{\OLId{latin-ordinary}{n}}) \land \Obj S_{\sigma}(\num{\OLId{latin-ordinary}{m}},
\num{\OLId{latin-ordinary}{n}})) \lif {}\\
& \qquad (\Obj Q_{\OLId{latin-ordinary}{q}'}(\num{\OLId{latin-ordinary}{m}}', \num{\OLId{latin-ordinary}{n}}') \land
\Obj S_{\sigma'}(\num{\OLId{latin-ordinary}{m}}, \num{\OLId{latin-ordinary}{n}}') \land !A(\num{\OLId{latin-ordinary}{m}}, \num{\OLId{latin-ordinary}{n}})).
\intertext{ويفيد فرض الاستقراء أن $!T(M,w)\Entails !C(M,w,n)$؛ أي}
& \Obj Q_\OLId{latin-ordinary}{q}(\num{\OLId{latin-ordinary}{m}}, \num{\OLId{latin-ordinary}{n}}) \land \Obj S_{\sigma_\OLMathNumeral{0}}(\num{\OLMathNumeral{0}}, \num{\OLId{latin-ordinary}{n}})
\land \dots \land \Obj S_{\sigma_\OLId{latin-ordinary}{k}}(\num{\OLId{latin-ordinary}{k}}, \num{\OLId{latin-ordinary}{n}}) \land \\
& \qquad\lforall[\OLId{latin-ordinary}{x}][(\num{\OLId{latin-ordinary}{k}} < \OLId{latin-ordinary}{x} \lif \Obj S_\TMblank(\OLId{latin-ordinary}{x}, \num{\OLId{latin-ordinary}{n}}))]\\
\intertext{ولما كانت الخانة~$m$ بعد $n$ خطوة تحتوي~$\sigma$، كان
  المقترن المقابل $\Obj S_\sigma(\num{m},\num{n})$، ولذلك تستلزم هذه}
& \Obj Q_{\OLId{latin-ordinary}{q}}(\num{\OLId{latin-ordinary}{m}}, \num{\OLId{latin-ordinary}{n}}) \land \Obj S_{\sigma}(\num{\OLId{latin-ordinary}{m}},
   \num{\OLId{latin-ordinary}{n}}).
\intertext{فنحصل الآن على}
& \Obj Q_{\OLId{latin-ordinary}{q}'}(\num{\OLId{latin-ordinary}{m}}', \num{\OLId{latin-ordinary}{n}}') \land \Obj S_{\sigma'}(\num{\OLId{latin-ordinary}{m}},
  \num{\OLId{latin-ordinary}{n}}') \land {}\\
& \qquad\displaystyle\bigwedge_{\substack{\OLMathNumeral{0}\leq \OLId{latin-ordinary}{i}\leq \OLId{latin-ordinary}{k}\\\OLId{latin-ordinary}{i}\neq \OLId{latin-ordinary}{m}}}
  \Obj S_{\sigma_\OLId{latin-ordinary}{i}}(\num{\OLId{latin-ordinary}{i}}, \num{\OLId{latin-ordinary}{n}}') \land {}\\
& \qquad \lforall[\OLId{latin-ordinary}{x}][(\num{\OLId{latin-ordinary}{k}} < \OLId{latin-ordinary}{x} \lif \Obj S_\TMblank(\OLId{latin-ordinary}{x}, \num{\OLId{latin-ordinary}{n}}'))]
\end{align*}
% تصحيح مصدر (C271-01): قُيّد حفظ رموز الشريط بالخانات i≠m؛ فالخانة m تُكتب فيها sigma'.
وبيان ذلك أن السطر الأول مأخوذ بقياس الفصل من تالي الشرط
المتقدم. وأما مقترنات السطر الأوسط، مع استثناء
$\OLId{latin-looped}{S}_{\sigma_\OLId{latin-ordinary}{m}}(\num{\OLId{latin-ordinary}{m}},\num{\OLId{latin-ordinary}{n}}')$، فيؤخذ كل واحد منها من
نظيره في $!C(\OLId{latin-looped}{M},\OLId{latin-ordinary}{w},\OLId{latin-ordinary}{n})$ مع $!A(\num{\OLId{latin-ordinary}{m}},\num{\OLId{latin-ordinary}{n}})$.

فإن كان $\OLId{latin-ordinary}{m}<\OLId{latin-ordinary}{k}$، كان $!T(\OLId{latin-looped}{M},\OLId{latin-ordinary}{w})\Proves\num{\OLId{latin-ordinary}{m}}<\num{\OLId{latin-ordinary}{k}}$
(\olref[rep]{prop:mlessk})، ويعطي تعدي~$<$
$\lforall[\OLId{latin-ordinary}{x}][(\num{\OLId{latin-ordinary}{k}}<\OLId{latin-ordinary}{x}\lif\num{\OLId{latin-ordinary}{m}}<\OLId{latin-ordinary}{x})]$. وإن كان $\OLId{latin-ordinary}{m}=\OLId{latin-ordinary}{k}$،
ثبت $\lforall[\OLId{latin-ordinary}{x}][(\num{\OLId{latin-ordinary}{k}}<\OLId{latin-ordinary}{x}\lif\num{\OLId{latin-ordinary}{m}}<\OLId{latin-ordinary}{x})]$ بالمنطق وحده.
فيؤخذ السطر الأخير من المقترن الموافق في $!C(\OLId{latin-looped}{M},\OLId{latin-ordinary}{w},\OLId{latin-ordinary}{n})$، ومن
$\lforall[\OLId{latin-ordinary}{x}][(\num{\OLId{latin-ordinary}{k}}<\OLId{latin-ordinary}{x}\lif\num{\OLId{latin-ordinary}{m}}<\OLId{latin-ordinary}{x})]$، ومن
$!A(\num{\OLId{latin-ordinary}{m}},\num{\OLId{latin-ordinary}{n}})$. وفي الحالة $\OLId{latin-ordinary}{m}<\OLId{latin-ordinary}{k}$ يكون الحاصل هو
$!C(\OLId{latin-looped}{M},\OLId{latin-ordinary}{w},\OLId{latin-ordinary}{n}+\OLMathNumeral{1})$ بلا زيادة.

وأما إذا كان $\OLId{latin-ordinary}{m}=\OLId{latin-ordinary}{k}$، فقد بلغ الرأس عند تمام $\OLId{latin-ordinary}{n}+\OLMathNumeral{1}$ خطوة
الخانة~$\OLId{latin-ordinary}{k}+\OLMathNumeral{1}$، فصارت أقصى ما زاره يمينًا. ولهذا يزاد في
$!C(\OLId{latin-looped}{M},\OLId{latin-ordinary}{w},\OLId{latin-ordinary}{n}+\OLMathNumeral{1})$ المقترن $\Obj S_\TMblank(\num{\OLId{latin-ordinary}{k}}',\num{\OLId{latin-ordinary}{n}}')$،
ويكون آخر مقترناتها
$\lforall[\OLId{latin-ordinary}{x}][(\num{\OLId{latin-ordinary}{k}}'<\OLId{latin-ordinary}{x}\lif\Obj S_\TMblank(\OLId{latin-ordinary}{x},\num{\OLId{latin-ordinary}{n}}'))]$.
فيلزمنا إثبات استلزام هاتين ال!!{sentence}s أيضًا.

قد حصل لنا $\lforall[\OLId{latin-ordinary}{x}][(\num{\OLId{latin-ordinary}{k}}<\OLId{latin-ordinary}{x}\lif\Obj S_\TMblank(\OLId{latin-ordinary}{x},
  \num{\OLId{latin-ordinary}{n}}'))]$؛ ومن تخصيصها
$\num{\OLId{latin-ordinary}{k}}<\num{\OLId{latin-ordinary}{k}}'\lif\Obj S_\TMblank(\num{\OLId{latin-ordinary}{k}}',\num{\OLId{latin-ordinary}{n}}')$.
والبديهية $\lforall[\OLId{latin-ordinary}{x}][\OLId{latin-ordinary}{x}<\OLId{latin-ordinary}{x}']$ تعطي $\num{\OLId{latin-ordinary}{k}}<\num{\OLId{latin-ordinary}{k}}'$،
فيتم بقياس الفصل $\Obj S_\TMblank(\num{\OLId{latin-ordinary}{k}}',\num{\OLId{latin-ordinary}{n}}')$.

ثم إن $!T(\OLId{latin-looped}{M},\OLId{latin-ordinary}{w})\Proves\num{\OLId{latin-ordinary}{k}}<\num{\OLId{latin-ordinary}{k}}'$، فبضم بديهية تعدي
$<$ إلى ما تقدم نحصل على
$\lforall[\OLId{latin-ordinary}{x}][(\num{\OLId{latin-ordinary}{k}}'<\OLId{latin-ordinary}{x}\lif\Obj S_\TMblank(\OLId{latin-ordinary}{x},\num{\OLId{latin-ordinary}{n}}'))]$.
وتفصيل هذا الاستنتاج متروك للتمرين.

\item وأما التعليمة من الصورة~\olref{left}، فبحسب
  \olref[rep]{defn:tm-descr}\olref[rep]{rep-left}،
\begin{align*}
& \lforall[\OLId{latin-ordinary}{x}][\lforall[\OLId{latin-ordinary}{y}][((\Obj Q_{\OLId{latin-ordinary}{q}}(\OLId{latin-ordinary}{x}', \OLId{latin-ordinary}{y}) \land \Obj
    S_{\sigma}(\OLId{latin-ordinary}{x}', \OLId{latin-ordinary}{y})) \lif {}]]\\
& \qquad (\Obj Q_{\OLId{latin-ordinary}{q}'}(\OLId{latin-ordinary}{x}, \OLId{latin-ordinary}{y}') \land \Obj
  S_{\sigma'}(\OLId{latin-ordinary}{x}', \OLId{latin-ordinary}{y}') \land !A(\OLId{latin-ordinary}{x}', \OLId{latin-ordinary}{y}))) \land {}\\
& \lforall[\OLId{latin-ordinary}{y}][((\Obj Q_{\OLId{latin-ordinary}{q}}(\num{\OLMathNumeral{0}}, \OLId{latin-ordinary}{y}) \land \Obj S_{\sigma}(\num{\OLMathNumeral{0}},
    \OLId{latin-ordinary}{y})) \lif {}]\\
& \qquad (\Obj Q_{\OLId{latin-ordinary}{q}'}(\num{\OLMathNumeral{0}}, \OLId{latin-ordinary}{y}') \land \Obj S_{\sigma'}(\num{\OLMathNumeral{0}}, \OLId{latin-ordinary}{y}')
  \land !A(\num{\OLMathNumeral{0}}, \OLId{latin-ordinary}{y})))
\intertext{مقترنة في $!T(M,w)$. فإذا كان $m>0$، فليكن $l=m-1$
  (أي $m=l+1$). ويستلزم المقترن الأول من ال!!{sentence} السابقة ما يأتي:}
& (\Obj Q_{\OLId{latin-ordinary}{q}}(\num{\OLId{latin-ordinary}{l}}', \num{\OLId{latin-ordinary}{n}}) \land \Obj S_{\sigma}(\num{\OLId{latin-ordinary}{l}}', \num{\OLId{latin-ordinary}{n}}))
\lif {} \\
& \qquad (\Obj Q_{\OLId{latin-ordinary}{q}'}(\num{\OLId{latin-ordinary}{l}}, \num{\OLId{latin-ordinary}{n}}') \land \Obj S_{\sigma'}(\num{\OLId{latin-ordinary}{l}}',
\num{\OLId{latin-ordinary}{n}}') \land !A(\num{\OLId{latin-ordinary}{l}}', \num{\OLId{latin-ordinary}{n}})).
\intertext{وإلا فليكن $l=m=0$، وانظر إلى ال!!{sentence} الآتية التي
  يستلزمها المقترن الثاني:}
& ((\Obj Q_{\OLId{latin-ordinary}{q}}(\num{\OLMathNumeral{0}}, \num{\OLId{latin-ordinary}{n}}) \land \Obj S_{\sigma}(\num{\OLMathNumeral{0}}, \num{\OLId{latin-ordinary}{n}})) \lif {}\\
& \qquad   (\Obj Q_{\OLId{latin-ordinary}{q}'}(\num{\OLMathNumeral{0}}, \num{\OLId{latin-ordinary}{n}}') \land \Obj S_{\sigma'}(\num{\OLMathNumeral{0}}, \num{\OLId{latin-ordinary}{n}}') \land
!A(\num{\OLMathNumeral{0}}, \num{\OLId{latin-ordinary}{n}}))).
\intertext{وتستلزم كل من العبارتين}
&  \Obj Q_{\OLId{latin-ordinary}{q}'}(\num{\OLId{latin-ordinary}{l}}, \num{\OLId{latin-ordinary}{n}}') \land \Obj S_{\sigma'}(\num{\OLId{latin-ordinary}{m}},
  \num{\OLId{latin-ordinary}{n}}') \land {}\\
&\qquad \displaystyle\bigwedge_{\substack{\OLMathNumeral{0}\leq \OLId{latin-ordinary}{i}\leq \OLId{latin-ordinary}{k}\\\OLId{latin-ordinary}{i}\neq \OLId{latin-ordinary}{m}}}
  \Obj S_{\sigma_\OLId{latin-ordinary}{i}}(\num{\OLId{latin-ordinary}{i}}, \num{\OLId{latin-ordinary}{n}}') \land {} \\
& \qquad  \lforall[\OLId{latin-ordinary}{x}][(\num{\OLId{latin-ordinary}{k}} < \OLId{latin-ordinary}{x}
  \lif \Obj S_\TMblank(\OLId{latin-ordinary}{x}, \num{\OLId{latin-ordinary}{n}}'))]
\end{align*}
% تصحيح مصدر (C271-02): رُبط فرع أقصى اليسار بالحالتين q وq' في التعليمة محل البحث.
على الوجه المتقدم. ففي الحالة الأولى
$\num{\OLId{latin-ordinary}{l}}'\ident\num{\OLId{latin-ordinary}{l}+\OLMathNumeral{1}}\ident\num{\OLId{latin-ordinary}{m}}$، وفي الثانية
$\num{\OLId{latin-ordinary}{l}}\ident\num{\OLMathNumeral{0}}$. والحاصل هو $!C(\OLId{latin-looped}{M},\OLId{latin-ordinary}{w},\OLId{latin-ordinary}{n}+\OLMathNumeral{1})$ نفسه.

\item وأما الحالة~\olref{stay} فبرهانها تمرين.
\end{enumerate}
فقد ثبت لكل~$\OLId{latin-ordinary}{n}$ إلى أول تهيئة توقف أن
$!T(\OLId{latin-looped}{M},\OLId{latin-ordinary}{w})\Entails !C(\OLId{latin-looped}{M},\OLId{latin-ordinary}{w},\OLId{latin-ordinary}{n})$.
\end{proof}

\begin{prob}
تمم حجة الحالة~\olref[tur][und][ver]{stay} من برهان
\olref[tur][und][ver]{lem:config}.
\end{prob}

\begin{prob}
أنشئ !!a{derivation} نتيجته $\Obj S_{\sigma_\OLId{latin-ordinary}{i}}(\num{\OLId{latin-ordinary}{i}},\num{\OLId{latin-ordinary}{n}}')$، ومقدماته
$\Obj S_{\sigma_\OLId{latin-ordinary}{i}}(\num{\OLId{latin-ordinary}{i}},\num{\OLId{latin-ordinary}{n}})$ و$!A(\OLId{latin-ordinary}{m},\OLId{latin-ordinary}{n})$ (على فرض
$\OLId{latin-ordinary}{i}\neq \OLId{latin-ordinary}{m}$؛ أي إما $\OLId{latin-ordinary}{i}<\OLId{latin-ordinary}{m}$ وإما $\OLId{latin-ordinary}{m}<\OLId{latin-ordinary}{i}$).
\end{prob}

\begin{prob}
أنشئ !!a{derivation} نتيجته
$\lforall[\OLId{latin-ordinary}{x}][(\num{\OLId{latin-ordinary}{k}}'<\OLId{latin-ordinary}{x}\lif\Obj S_\TMblank(\OLId{latin-ordinary}{x},\num{\OLId{latin-ordinary}{n}}'))]$، ومقدماته
$\lforall[\OLId{latin-ordinary}{x}][(\num{\OLId{latin-ordinary}{k}}<\OLId{latin-ordinary}{x}\lif\Obj S_\TMblank(\OLId{latin-ordinary}{x},\num{\OLId{latin-ordinary}{n}}'))]$، و
$\lforall[\OLId{latin-ordinary}{x}][\OLId{latin-ordinary}{x}<\OLId{latin-ordinary}{x}']$، و
$\lforall[\OLId{latin-ordinary}{x}][\lforall[\OLId{latin-ordinary}{y}][\lforall[\OLId{latin-ordinary}{z}][((\OLId{latin-ordinary}{x}<\OLId{latin-ordinary}{y}\land \OLId{latin-ordinary}{y}<\OLId{latin-ordinary}{z})\lif \OLId{latin-ordinary}{x}<\OLId{latin-ordinary}{z})]]]$.)
\end{prob}

\begin{lem}
\ollabel{lem:valid-if-halt}
توقف $\OLId{latin-looped}{M}$ عند الدخل~$\OLId{latin-ordinary}{w}$ يقتضي أن تكون $!T(\OLId{latin-looped}{M},\OLId{latin-ordinary}{w})\lif !E(\OLId{latin-looped}{M},\OLId{latin-ordinary}{w})$ صادقة
منطقيًا.
\end{lem}

\begin{proof}
بمقتضى \olref{lem:config}، لكل زمن~$\OLId{latin-ordinary}{n}$ إلى أول تهيئة توقف،
يكون الوصف~$!C(\OLId{latin-looped}{M},\OLId{latin-ordinary}{w},\OLId{latin-ordinary}{n})$ لتهيئة $\OLId{latin-looped}{M}$ في الزمن~$\OLId{latin-ordinary}{n}$ لازمًا
من~$!T(\OLId{latin-looped}{M},\OLId{latin-ordinary}{w})$. ولتتوقف $\OLId{latin-looped}{M}$ بعد $\OLId{latin-ordinary}{k}$ خطوات، ورأسها على
الخانة~$\OLId{latin-ordinary}{m}$ لبعض $\OLId{latin-ordinary}{m}\in\Nat$. عندئذ تصف $!C(\OLId{latin-looped}{M},\OLId{latin-ordinary}{w},\OLId{latin-ordinary}{k})$ تهيئة
توقف~$\OLId{latin-looped}{M}$؛ فهي تتضمن المقترنين
$\Obj Q_\OLId{latin-ordinary}{q}(\num{\OLId{latin-ordinary}{m}},\num{\OLId{latin-ordinary}{k}})$ و$\Obj S_\sigma(\num{\OLId{latin-ordinary}{m}},\num{\OLId{latin-ordinary}{k}})$،
مع كون $\delta(\OLId{latin-ordinary}{q},\sigma)$ غير معرّفة. فتعطي
\olref{lem:halt-config-implies-halt}
$!C(\OLId{latin-looped}{M},\OLId{latin-ordinary}{w},\OLId{latin-ordinary}{k})\Entails !E(\OLId{latin-looped}{M},\OLId{latin-ordinary}{w})$. وقد ثبت
$!T(\OLId{latin-looped}{M},\OLId{latin-ordinary}{w})\Entails !C(\OLId{latin-looped}{M},\OLId{latin-ordinary}{w},\OLId{latin-ordinary}{k})$، فينتج
$!T(\OLId{latin-looped}{M},\OLId{latin-ordinary}{w})\Entails !E(\OLId{latin-looped}{M},\OLId{latin-ordinary}{w})$، وتكون
$!T(\OLId{latin-looped}{M},\OLId{latin-ordinary}{w})\lif !E(\OLId{latin-looped}{M},\OLId{latin-ordinary}{w})$ صادقة منطقيًا.
\end{proof}

\begin{explain}
وبقي لإتمام التحقق عكس هذه النتيجة: من الصدق المنطقي لـ
$!T(\OLId{latin-looped}{M},\OLId{latin-ordinary}{w})\lif !E(\OLId{latin-looped}{M},\OLId{latin-ordinary}{w})$ يلزم أن $\OLId{latin-looped}{M}$ تتوقف فعلًا عند
الدخل~$\OLId{latin-ordinary}{w}$.
\end{explain}

\begin{lem}
\ollabel{lem:halt-if-valid}
من $\Entails !T(\OLId{latin-looped}{M},\OLId{latin-ordinary}{w})\lif !E(\OLId{latin-looped}{M},\OLId{latin-ordinary}{w})$ يلزم توقف $\OLId{latin-looped}{M}$ عند
الدخل~$\OLId{latin-ordinary}{w}$.
\end{lem}

\begin{proof}
لنأخذ ال!!{structure}~$\Struct M$ للغة $\Lang L_\OLId{latin-looped}{M}$ ذات
المجال $\Nat$؛ نفسر فيها $\Obj 0$ بالعدد~$\OLMathNumeral{0}$،
و$\prime$ بدالة الخلف، و$<$ بعلاقة الأصغر من. ونجعل تفسير
المحمولين $\Obj Q_\OLId{latin-ordinary}{q}$ و~$\Obj S_\sigma$ على الوجه الآتي:
\begin{align*}
  \Assign{\Obj Q_\OLId{latin-ordinary}{q}}{\OLId{latin-looped}{M}} & =
\Setabs{\tuple{\OLId{latin-ordinary}{m}, \OLId{latin-ordinary}{n}}}{\begin{array}{ll}\text{بعد بدء $M$ عند~$w$ ومرور~$n$ خطوة،}\\ \text{تكون في الحالة $q$
  وتمسح الخانة~$m$}\end{array}} \\
\Assign{\Obj S_\sigma}{\OLId{latin-looped}{M}} & = \Setabs{\tuple{\OLId{latin-ordinary}{m}, \OLId{latin-ordinary}{n}}}{\begin{array}{ll}
\text{بعد بدء $M$ عند~$w$ ومرور $n$ خطوة،}\\ \text{تحتوي الخانة~$m$ من شريط $M$
  الرمز~$\sigma$}\end{array}}
\end{align*}
فقد بنينا ال!!{structure}~$\Struct{M}$ لتطابق، خطوة بعد خطوة،
ما تفعله $\OLId{latin-looped}{M}$ حقًا عند الدخل~$\OLId{latin-ordinary}{w}$. ومن البناء يتضح
$\Sat{\OLId{latin-looped}{M}}{!T(\OLId{latin-looped}{M},\OLId{latin-ordinary}{w})}$. فإن كان $\Entails !T(\OLId{latin-looped}{M},\OLId{latin-ordinary}{w})\lif !E(\OLId{latin-looped}{M},\OLId{latin-ordinary}{w})$،
لزم $\Sat{\OLId{latin-looped}{M}}{!E(\OLId{latin-looped}{M},\OLId{latin-ordinary}{w})}$، أي:
\[
\Sat{\OLId{latin-looped}{M}}{\lexists[\OLId{latin-ordinary}{x}][\lexists[\OLId{latin-ordinary}{y}][(\bigvee_{\tuple{\OLId{latin-ordinary}{q}, \sigma} \in
      \OLId{latin-looped}{X}}(\Obj Q_\OLId{latin-ordinary}{q}(\OLId{latin-ordinary}{x}, \OLId{latin-ordinary}{y}) \land \Obj S_\sigma(\OLId{latin-ordinary}{x}, \OLId{latin-ordinary}{y})))]]}.
\]
وحيث إن $\Domain{\OLId{latin-looped}{M}}=\Nat$، فلا بد من $\OLId{latin-ordinary}{m},\OLId{latin-ordinary}{n}\in\Nat$ يتحقق لهما
$\Sat{\OLId{latin-looped}{M}}{\Obj Q_\OLId{latin-ordinary}{q}(\num{\OLId{latin-ordinary}{m}},\num{\OLId{latin-ordinary}{n}})\land
\Obj S_\sigma(\num{\OLId{latin-ordinary}{m}},\num{\OLId{latin-ordinary}{n}})}$ لبعض~$\OLId{latin-ordinary}{q}$ و~$\sigma$،
وتكون $\delta(\OLId{latin-ordinary}{q},\sigma)$ غير معرّفة. وهذا، بتعريف~$\Struct M$،
معناه أن $\OLId{latin-looped}{M}$ تبدأ عند الدخل~$\OLId{latin-ordinary}{w}$، ثم تكون بعد~$\OLId{latin-ordinary}{n}$ خطوة في
الحالة~$\OLId{latin-ordinary}{q}$ قارئةً الرمز~$\sigma$، ولا انتقال لها؛ أي إن~$\OLId{latin-looped}{M}$
قد توقفت.
\end{proof}

\end{document}
